\documentclass[11pt]{article}

\usepackage[T1]{fontenc}
\usepackage{lmodern}
\usepackage{microtype}
\usepackage[a4paper,margin=27mm]{geometry}
\usepackage{amsmath,amssymb,amsthm,mathtools}
\usepackage{booktabs}
\usepackage{graphicx}
\usepackage{xcolor}
\usepackage{doi}
\usepackage{hyperref}

\hypersetup{
  colorlinks=true,
  linkcolor=blue!45!black,
  citecolor=blue!45!black,
  urlcolor=blue!55!black,
  pdftitle={An exact qubit-memory counterexample to an alpha-Doeblin hockey-stick bound},
  pdfauthor={Xiaojian Zeng}
}

\newtheorem{theorem}{Theorem}
\newtheorem{lemma}{Lemma}
\newcommand{\cN}{\mathcal N}
\newcommand{\id}{\operatorname{id}}
\newcommand{\Tr}{\operatorname{Tr}}
\newcommand{\ket}[1]{\lvert #1\rangle}
\newcommand{\bra}[1]{\langle #1\rvert}
\newcommand{\dyad}[2]{\lvert #1\rangle\!\langle #2\rvert}
\newcommand{\psdle}{\preceq}
\newcommand{\psdge}{\succeq}

\title{An exact qubit-memory counterexample to an
$\alpha$-Doeblin hockey-stick bound}
\author{Xiaojian Zeng\\
\small PBC School of Finance, Tsinghua University, Beijing 100190, China\\
\small \texttt{xiaojianzeng@foxmail.com}}
\date{Draft of 13 September 2026}

\begin{document}

\maketitle

\begin{abstract}
The complete contraction bound for quantum hockey-stick divergence is known
with the positive Doeblin coefficient $\alpha_+$.  A published open question
asks whether the stronger coefficient $\alpha$ can replace $\alpha_+$ when
the two inputs have the same reference marginal.  We give an exact negative
answer.  For one use of the qubit amplitude-damping channel with
$p=1$ and $\eta=1/4$, a noiseless qubit reference, bias $\gamma=2$, and the
entangled pair
$\ket{\psi_\pm}=(\ket{00}\pm10\ket{11})/\sqrt{101}$, we certify
$\alpha=1/4$, $\alpha_+=0$, and
\[
 E_2\bigl((\id\otimes\cN)(\rho)\Vert(\id\otimes\cN)(\sigma)\bigr)
 -(1-\alpha)E_2(\rho\Vert\sigma)
 =\frac{199+24\sqrt{41}-9\sqrt{1489}}{808}>0.
\]
All channel, state, marginal, semidefinite-program, and strict-sign
conditions are checked algebraically.  The witness disproves the universal
$\alpha$ replacement, without claiming a classification or the optimal
complete contraction coefficient.
\end{abstract}

\noindent\textbf{Keywords:} quantum channels; strong data processing;
hockey-stick divergence; Doeblin coefficient; quantum side information

\section{Introduction}

A contraction coefficient quantifies how much a channel can preserve the
distinguishability of two inputs.  Such coefficients strengthen ordinary
data processing and enter problems ranging from mixing to noisy quantum
hypothesis testing.  George, Hirche, Nuradha, and Wilde introduced several
quantum analogues of the classical Doeblin coefficient and related them to
efficient upper bounds on contraction coefficients~\cite{GeorgeEtAl2026}.
Two variants matter here.  The coefficient $\alpha$ optimizes over a
Hermitian operator, whereas $\alpha_+$ restricts that operator to be positive
semidefinite.  Consequently $\alpha_+\leq\alpha$, and a contraction factor
$1-\alpha$ is potentially strictly stronger than $1-\alpha_+$.

For the hockey-stick divergence, the known complete bound uses
$1-\alpha_+$.  Open Question~1 of Ref.~\cite{GeorgeEtAl2026} asks whether it
remains valid with $1-\alpha$.  The question is nontrivial precisely because
an optimizer for $\alpha$ need not be positive semidefinite and therefore
need not define a physical replacer channel.  Follow-up work on conditional
contraction retains an $\alpha_+$ bound~\cite{HircheEtAl2026Conditional},
while nonlinear hockey-stick inequalities obtained under an additional noise
criterion concern a different, non-reference setting~\cite{NuradhaEtAl2025Nonlinear}.

The main contribution is stated and proved with exact arithmetic.  A single
amplitude-damping channel has $\alpha=1/4$ but $\alpha_+=0$.  Two entangled
two-qubit inputs with identical reference marginals then violate the proposed
$1-\alpha$ bound at $\gamma=2$.  The proof supplies primal and dual
semidefinite-program certificates, exact input and output spectra, an
attaining measurement effect, and an integer/radical certificate for the
strict sign.  Because the proposed statement is universal, this one witness
settles it negatively.  It does not determine the best replacement factor,
describe all violations, or assert failure for every $\gamma>1$.

\section{Definitions and the proposed complete bound}

Let $\cN:\mathcal L(A)\to\mathcal L(B)$ be a completely positive,
trace-preserving map.  We use the unnormalized Choi convention
\begin{equation}
 J_{\cN}=\sum_{i,j}\dyad{i}{j}_A\otimes\cN(\dyad{i}{j}),
 \qquad J_{\cN}\psdge0,\qquad \Tr_BJ_{\cN}=I_A .
 \label{eq:choi}
\end{equation}
The two coefficients are
\begin{align}
 \alpha(\cN)
 &=\max_{H=H^\dagger}\bigl\{\Tr H:I_A\otimes H\psdle J_{\cN}\bigr\},
 \label{eq:alpha}\\
 \alpha_+(\cN)
 &=\max_{H\psdge0}\bigl\{\Tr H:I_A\otimes H\psdle J_{\cN}\bigr\}.
 \label{eq:alphaplus}
\end{align}
For density operators $\rho$ and $\sigma$ and $\gamma>1$, the hockey-stick
divergence is
\begin{equation}
 E_\gamma(\rho\Vert\sigma)
 =\Tr(\rho-\gamma\sigma)_+
 =\max_{0\psdle M\psdle I}\Tr M(\rho-\gamma\sigma).
 \label{eq:hockey}
\end{equation}

The established complete inequality says that, for states on $R\otimes A$
with $\rho_R=\sigma_R$,
\begin{equation}
 E_\gamma\!\left((\id_R\otimes\cN)(\rho)\middle\|
                  (\id_R\otimes\cN)(\sigma)\right)
 \leq[1-\alpha_+(\cN)]E_\gamma(\rho\Vert\sigma)
 \label{eq:known}
\end{equation}
for every $\gamma\geq1$~\cite[Proposition~31]{GeorgeEtAl2026}.  The proposed
strengthening replaces $\alpha_+$ by $\alpha$:
\begin{equation}
 E_\gamma\!\left((\id_R\otimes\cN)(\rho)\middle\|
                  (\id_R\otimes\cN)(\sigma)\right)
 \stackrel{?}{\leq}[1-\alpha(\cN)]E_\gamma(\rho\Vert\sigma).
 \label{eq:proposed}
\end{equation}
No ratio is taken in Eq.~\eqref{eq:proposed}; inputs with zero divergence
remain harmless by ordinary data processing.

For priors $1/(1+\gamma)$ on $\rho$ and $\gamma/(1+\gamma)$ on $\sigma$,
the Helstrom--Holevo binary error probability~\cite{Helstrom1967,Holevo1973}
is
\begin{equation}
 P_{\rm err}(\rho,\sigma)=\frac{1-E_\gamma(\rho\Vert\sigma)}{1+\gamma}.
 \label{eq:perr}
\end{equation}
Thus Eq.~\eqref{eq:proposed} would lower-bound the post-channel error in a
biased discrimination experiment assisted by a noiseless quantum memory.

\section{A channel with \texorpdfstring{$\alpha_+<\alpha$}{alpha-plus less than alpha}}

Consider the generalized amplitude-damping family of
Ref.~\cite[Lemma~10]{GeorgeEtAl2026} at $p=1$ and $\eta=1/4$.  Its two
nonzero Kraus operators are
\begin{equation}
 K_0=\begin{pmatrix}1&0\\0&1/2\end{pmatrix},\qquad
 K_1=\begin{pmatrix}0&\sqrt3/2\\0&0\end{pmatrix},
 \qquad \sum_aK_a^\dagger K_a=I.
 \label{eq:kraus}
\end{equation}
In the ordered Choi basis
$(\ket{0_A0_B},\ket{0_A1_B},\ket{1_A0_B},\ket{1_A1_B})$,
\begin{equation}
 J=\begin{pmatrix}
 1&0&0&1/2\\
 0&0&0&0\\
 0&0&3/4&0\\
 1/2&0&0&1/4
 \end{pmatrix}.
 \label{eq:J}
\end{equation}
It is positive semidefinite, with spectrum
$\{0,0,3/4,5/4\}$, and $\Tr_BJ=I_A$.

\begin{lemma}
For the channel in Eq.~\eqref{eq:kraus},
\begin{equation}
 \alpha(\cN)=\frac14,\qquad \alpha_+(\cN)=0.
 \label{eq:coefficients}
\end{equation}
\end{lemma}

\begin{proof}
For $\alpha$, choose
\begin{equation}
 H=\operatorname{diag}(1/2,-1/4),\qquad
 Y=\begin{pmatrix}
 1&0&0&-1\\0&0&0&0\\0&0&0&0\\-1&0&0&1
 \end{pmatrix}.
 \label{eq:alpha-cert}
\end{equation}
The primal slack $J-I\otimes H$ has spectrum
$\{0,1/4,1/4,1\}$, so $H$ is feasible and gives $\Tr H=1/4$.
Moreover $Y\psdge0$, $\Tr_A Y=I_B$, and $\Tr(JY)=1/4$.
The matching primal and dual objectives certify $\alpha=1/4$.  This also
reproduces the family formula
$\alpha(\cN_{p,\eta})=(1-\sqrt\eta)^2$ at the selected point.

For $\alpha_+$, the zero matrix is primal feasible.  Let
\begin{equation}
 \ket{\kappa}=\ket{00}-2\ket{11},\qquad
 Y_+=\ket{\kappa}\!\bra{\kappa}.
\end{equation}
Then $J\ket{\kappa}=0$ and
$\Tr_A Y_+=\operatorname{diag}(1,4)\psdge I$.  For every positive
primal-feasible $H_+$,
\begin{equation}
 \Tr H_+
 \leq\Tr[(\Tr_A Y_+)H_+]
 =\Tr[Y_+(I\otimes H_+)]
 \leq\Tr(Y_+J)=0.
\end{equation}
Together with the zero primal point, this proves $\alpha_+=0$.
\end{proof}

\section{Exact counterexample}

All state matrices below use the $R$-major order
$(\ket{0_R0_A},\ket{0_R1_A},\ket{1_R0_A},\ket{1_R1_A})$ before the channel,
and the corresponding $R,B$ order afterwards.

\begin{theorem}
Equation~\eqref{eq:proposed} is false even when
$A=B=R=\mathbb C^2$.  A counterexample is given by the channel in
Eq.~\eqref{eq:kraus}, $\gamma=2$, and
\begin{equation}
 \ket{\psi_\pm}=\frac{\ket{00}\pm10\ket{11}}{\sqrt{101}},
 \qquad
 \rho=\ket{\psi_+}\!\bra{\psi_+},\quad
 \sigma=\ket{\psi_-}\!\bra{\psi_-}.
 \label{eq:states}
\end{equation}
\end{theorem}

\begin{proof}
The two operators in Eq.~\eqref{eq:states} are rank-one density operators and
have the identical reference marginal
\begin{equation}
 \rho_R=\sigma_R=\operatorname{diag}(1/101,100/101).
 \label{eq:marginal}
\end{equation}
The eigenvalues of either partial transpose include $-10/101$, so the inputs
are entangled and cannot reduce to a common classical-reference construction.

The only nonzero block of $\rho-2\sigma$, on
$\operatorname{span}\{\ket{00},\ket{11}\}$, is
\begin{equation}
 \frac1{101}\begin{pmatrix}-1&30\\30&-100\end{pmatrix}.
\end{equation}
Its trace and determinant are $-1$ and $-800/101^2$, respectively, so it has
exactly one positive eigenvalue.  Direct diagonalization gives
\begin{equation}
 E_2(\rho\Vert\sigma)=\frac{3\sqrt{1489}-101}{202}.
 \label{eq:inputE}
\end{equation}

The channel action is
\begin{equation}
 \cN\!\begin{pmatrix}a&b\\c&d\end{pmatrix}
 =\begin{pmatrix}a+3d/4&b/2\\c/2&d/4\end{pmatrix}.
 \label{eq:action}
\end{equation}
Applying it to the $A$ subsystem gives
\begin{align}
 \rho'&=\frac1{101}\begin{pmatrix}
 1&0&0&5\\0&0&0&0\\0&0&75&0\\5&0&0&25
 \end{pmatrix},
 &
 \sigma'&=\frac1{101}\begin{pmatrix}
 1&0&0&-5\\0&0&0&0\\0&0&75&0\\-5&0&0&25
 \end{pmatrix}.
 \label{eq:outputs}
\end{align}
The active block of $\rho'-2\sigma'$ is
$101^{-1}\bigl(\begin{smallmatrix}-1&15\\15&-25\end{smallmatrix}\bigr)$;
the remaining eigenvalues are $0$ and $-75/101$.  Hence
\begin{equation}
 E_2(\rho'\Vert\sigma')=\frac{3\sqrt{41}-13}{101}.
 \label{eq:outputE}
\end{equation}
For completeness, the positive eigenspace is attained by the effect whose
active two-dimensional block is
\begin{equation}
 M_{\rm act}=\frac12\left[I+\frac{5\sigma_x+4\sigma_z}{\sqrt{41}}\right].
 \label{eq:effect}
\end{equation}
Indeed $M_{\rm act}^2=M_{\rm act}$ because $5^2+4^2=41$; embedding it by
zero on the other two basis vectors gives $0\psdle M\psdle I$ and reproduces
Eq.~\eqref{eq:outputE} through Eq.~\eqref{eq:hockey}.

Using Eq.~\eqref{eq:coefficients}, the excess over the proposed right-hand
side is
\begin{equation}
 \Delta
 =\frac{3\sqrt{41}-13}{101}
  -\frac34\frac{3\sqrt{1489}-101}{202}
 =\frac{199+24\sqrt{41}-9\sqrt{1489}}{808}.
 \label{eq:margin}
\end{equation}
It remains to certify strictness without a floating-point comparison.  Both
sides of $199+24\sqrt{41}>9\sqrt{1489}$ are positive, and their squared
difference is
\begin{equation}
 -57392+9552\sqrt{41}
 >-57392+9552\frac{32}{5}
 =\frac{18704}{5}>0,
\end{equation}
where $\sqrt{41}>32/5$ follows by squaring.  Therefore $\Delta>0$, which
contradicts Eq.~\eqref{eq:proposed}.
\end{proof}

\section{Discrimination meaning and comparison with known bounds}

Figure~\ref{fig:comparison} displays decimal renderings of the exact values
only for visual comparison.  Ordinary data processing remains intact:
$E_2(\rho'\Vert\sigma')\approx0.061479$ is below the input value
$0.073083$.  The known $\alpha_+$ inequality~\eqref{eq:known} also remains
valid, since $\alpha_+=0$.  What fails is the stronger proposed ceiling
$(1-\alpha)E_2\approx0.054812$.

\begin{figure}[t]
 \centering
 \includegraphics[width=0.88\linewidth]{figures/witness_comparison.pdf}
 \caption{Exact witness values, plotted in decimal form.  The actual output
 divergence exceeds the proposed $\alpha$-based ceiling by the exact positive
 quantity $\Delta$ in Eq.~\eqref{eq:margin}.  The plot is generated directly
 from the symbolic expressions; no numerical fit or optimization enters the
 proof.}
 \label{fig:comparison}
\end{figure}

For priors $1/3$ on $\rho$ and $2/3$ on $\sigma$, Eq.~\eqref{eq:perr} gives
an actual post-channel error of approximately $0.312840$.  If
Eq.~\eqref{eq:proposed} held, it would require an error of at least
$0.315063$.  The proposed guarantee therefore overstates the unavoidable
loss of biased discrimination power by exactly $\Delta/3$ for this
memory-assisted experiment.

The role of the quantum reference is essential to the example.  Without a
reference system, the ordinary hockey-stick contraction coefficient is
bounded by the trace-distance coefficient and hence by
$1-\alpha$~\cite[Section~6.1]{GeorgeEtAl2026}; a common classical reference
basis likewise reduces the problem blockwise~\cite[Appendix~A]{HircheEtAl2026Conditional}.
The negative partial-transpose eigenvalue in Eq.~\eqref{eq:marginal} rules out
those reductions.  The strict gap $\alpha_+<\alpha$ also rules out the region
already covered by equality of the two Doeblin coefficients.

Among the directly related works checked through 13 September 2026, the
conditional framework of Ref.~\cite{HircheEtAl2026Conditional} continues to
use $\alpha_+$ and explicitly discusses the obstruction posed by a quantum
reference at general $\gamma$.  The nonlinear bounds of
Ref.~\cite{NuradhaEtAl2025Nonlinear} assume an additional channel-noise
condition and optimize without the same-reference complete setting.  These
comparisons are bounded rather than an exhaustive priority search; the result
is therefore presented without a priority claim.

\section{Scope and reproducibility}

The theorem uses one exact, finite-dimensional channel instance and one exact
pair of input states.  It is sufficient to disprove both the qubit-universal
statement and the arbitrary-dimensional universal proposal that contains it.
It does not characterize the optimal complete contraction coefficient, all
violating channels or inputs, the full range of $\gamma$, or robustness to
experimental imperfections.  No perturbative, asymptotic, continuum, or
unit-dependent statement is involved.

The ancillary package contains two independent exact implementations.  One
constructs the Kraus action, Choi matrix, semidefinite-program certificates,
states, spectra, and effect using symbolic rational and radical arithmetic.
The other begins from the rational state matrices and the block action in
Eq.~\eqref{eq:action}, imports no result from the first checker, and repeats
the decisive characteristic-polynomial and sign checks.  A separate script
generates Fig.~\ref{fig:comparison} and its tabulated values from the exact
expressions.  The environment and commands are specified in
\texttt{REPRODUCIBILITY.md}.

\section*{Data and code availability}

All witness matrices are included in the ancillary JSON data, and all exact
verification and figure-generation scripts are included in the local source
package.  No downloaded third-party paper is part of that package.  A public
archival location and license for the code must be selected by the prospective
author before dissemination or submission.

\section*{Author contributions and AI-use disclosure}

This is a provisional local draft for author review.  Xiaojian Zeng is the
prospective corresponding author.  OpenAI Codex language models were used for
idea development, the exact calculations and code, a bounded bibliographic
check, figure production, and drafting and editing the manuscript.  The models
are not authors.  Before any dissemination or submission, the prospective
author must independently inspect the sources, calculations, code, claims,
text, and rights; make any needed corrections; and explicitly accept full
responsibility.  That human review and approval are not represented as
complete here.

\section*{Competing interests}

The prospective author has not yet supplied a competing-interest declaration;
it must be completed before submission.

\bibliographystyle{unsrturl}
\bibliography{manuscript/references}

\end{document}
